% !PRINTED 295
\weberpage{295}
%<<CONT
\gesperrt{kann diese Gleichung nur dann durch successive Adjunction von Wurzeln
cyklischer Gleichungen reducibel werden, wenn sie imprimitiv ist.}

Wir können über die Art und Weise der Reduction, ihre Möglichkeit vorausgesetzt,
noch einiges Nähere anführen. Ist $P_\nu$ die erste Gruppe der Reihe (8), die
Normaltheiler von $P$ ist, so ist nach dem eben angeführten Satze III, §.\,7 bei
richtiger Anordnung der Compositionsreihe:
\weq{9}{j_\lambda = j_{\lambda+1} \cdots = j_{\nu+1},}
also alle gleich derselben Primzahl, und die Resolvente $\chi = 0$, durch die die
Gruppe $P$ auf $P_\nu$ reducirt wird, hat folglich eine metacyklische Gruppe.

Bezeichnen wir mit $A, B, \ldots, S$ die Systeme der Intransitivität der Gruppe
$P_\nu$, deren Anzahl $s$ sei, und setzen
\weq{10}{m = r\, s,}
so wird durch Adjunction einer Wurzel von $\chi = 0$ die Function $f(x)$ in $s$
Factoren $r^{\mathrm{ten}}$ Grades zerfallen.

Die $A, B, \ldots, S$ sind Systeme der Imprimitivität von $P$ (Bd. I, §.\,158, 2.).

Bezeichnen wir mit $Q$ die Gesammtheit aller Permutationen von $P$, die die
einzelnen Systeme $A, B, \ldots, S$ an ihrer Stelle lassen und nur die Elemente
der Systeme unter sich vertauschen, so ist $Q$, wie wir im §.\,158 des ersten
Bandes gesehen haben, ein Normaltheiler von $P$, und durch Adjunction der Wurzeln
einer Hülfsgleichung $s^{\mathrm{ten}}$ Grades $\varphi(y) = 0$ zerfällt $f(x)$ in
$s$ Factoren $r^{\mathrm{ten}}$ Grades, denen die einzelnen Systeme
$A, B, \ldots, S$ als Wurzeln angehören:
\[
f(x) = f(x, y_\alpha)\, f(x, y_\beta)\, f(x, y_\gamma) \ldots
\]
Da durch die Hülfsgleichung $\varphi(y) = 0$ die Gruppe $P$ auf $Q$ reducirt wird,
so ist die Gruppe von $\varphi(y) = 0$ isomorph mit $P/Q$. Andererseits ist aber
auch die intransitive Gruppe $P_\nu$ ein Theiler von $Q$, da die Systeme
$A, B, \ldots, S$ durch $P_\nu$ nicht verschoben werden, und wir können also den
Satz §.\,73, 3. auf die drei Gruppen $P, Q, P_\nu$ anwenden. Da die Indexreihe
von $P/P_\nu$ nach Voraussetzung aus lauter Primzahlen besteht, so gilt nach
jenem Satze dasselbe von $P/Q$. Diese Gruppe ist metacyklisch und also ist auch
die Hülfsgleichung $\varphi(y) = 0$ metacyklisch. Es folgt also der Satz:
%CONT>>
